\chapter*{Abstract}
Traffic Sign Recognition (TSR) is a critical component of Advanced Driver Assistance Systems (ADAS) and autonomous vehicle safety. This report presents a comprehensive investigation into a hybrid classification pipeline that leverages deep feature representations combined with Support Vector Machine (SVM) strategies. Utilizing a representative 10-class subset of the German Traffic Sign Recognition Benchmark (GTSRB), we evaluate two state-of-the-art Convolutional Neural Network (CNN) architectures, ResNet50 and InceptionV3, as automated feature extractors. The high-dimensional feature manifolds are analyzed using dimensionality reduction techniques, specifically Principal Component Analysis (PCA) and Uniform Manifold Approximation and Projection (UMAP), to assess class separability. Experimental results demonstrate that the InceptionV3-SVM hybrid architecture achieves a peak classification accuracy of \textbf{97.24\%}, outperforming the linear SGD-SVM baseline. Furthermore, we position our results against traditional handcrafted feature engineering methods, such as multi-scale HOG and Random Forests, documented in classic literature. Our findings confirm that pre-trained deep manifolds provide a robust and efficient alternative to labor-intensive feature engineering while maintaining high predictive performance in complex environmental conditions.

\vspace{0.5cm}
\noindent \textbf{Open Source:} All source code, modular experimental results, and interactive 3D visualizations are available at: \url{https://github.com/toando0212/lab-ml2}
